El presente informe aborda el estudio empírico de dos problemas clásicos en computación: el ordenamiento de arreglos unidimensionales y la multiplicación de matrices cuadradas. Se contrastan algoritmos de ordenamiento (\texttt{std::sort}, MergeSort, QuickSort y PatienceSort) junto a algoritmos matriciales (Naive iterativo y Strassen recursivo) sobre diferentes tamaños de entrada, distribuciones y dominios de valores. El propósito central no radica en demostrar analíticamente sus cotas asintóticas conocidas, sino en examinar cómo factores reales de la arquitectura moderna de computadores —tales como la jerarquía de memoria caché, la sobrecarga por asignación dinámica en el \textit{heap} y las optimizaciones a nivel de compilador— modulan, amplifican o contradicen las predicciones teóricas en la práctica.